\documentclass[12pt,a4paper]{article}

\usepackage{amsmath,amssymb,amsthm}
\usepackage{mathtools}
\usepackage{geometry}
\geometry{margin=2.5cm}
\usepackage{hyperref}
\usepackage{cleveref}
\usepackage{enumitem}
\usepackage{doi}

\newtheorem{theorem}{Theorem}[section]
\newtheorem{proposition}[theorem]{Proposition}
\newtheorem{lemma}[theorem]{Lemma}
\newtheorem{corollary}[theorem]{Corollary}
\theoremstyle{definition}
\newtheorem{definition}[theorem]{Definition}
\newtheorem{remark}[theorem]{Remark}

\newcommand{\Spi}{S_{\Pi}}
\newcommand{\Ag}{A_g}
\newcommand{\hk}{K}
\newcommand{\tr}{\operatorname{Tr}}
\newcommand{\vol}{\operatorname{vol}}
\newcommand{\dd}{\mathrm{d}}
\newcommand{\half}{\tfrac{1}{2}}

\begin{document}

  \title{Causal Propagation and Gravitational Waves from Projective Spectral Dynamics}

  \author{J\'{e}r\^{o}me Beau\\
  \small Independent Researcher, France\\
  \small \texttt{jerome.beau@cosmochrony.org}}

  \maketitle

  \abstract{
  We construct the Lorentzian completion of the spectral entropy framework
  in which the renormalized variation of a Laplace-type operator generates
  an infrared Einstein response.
  While the original formulation was elliptic and sufficient to derive
  the effective field equations and their thermodynamic interpretation,
  it did not address causal propagation.
  We define a Lorentzian spectral action via the real part of the
  hyperbolic determinant, equivalently through a Schwinger--Keldysh
  prescription, ensuring a real and retarded quadratic kernel.

  Linearization about Minkowski spacetime shows that, in the regime
  $R\ell_\chi^2 \ll 1$ and $\omega \ll \ell_\chi^{-1}$,
  the theory propagates exactly two transverse-traceless tensor modes of
  helicity $\pm2$.
  Scalar and vector sectors remain non-dynamical in the infrared, and
  the additional pole generated by higher-derivative corrections lies at
  $k^2 \sim \ell_\chi^{-2}$, beyond current observational reach.
  The polarization content therefore coincides with that of general relativity
  in the accessible domain.
  When the vector sector is sourced by a projected momentum density it obeys an elliptic gravitomagnetic
  constraint, fixing the frame-dragging potential $g_{0i}$ without an additional propagating degree of freedom,
  conditionally on the projected matter current.

  Beyond the leading Einstein sector, the ratio of Seeley--DeWitt
  coefficients fixes a universal $k^4$ correction to the dispersion relation,
  \[
    \omega^2 = c^2 k^2 - \gamma \ell_\chi^2 k^4 + \mathcal{O}(k^6),
  \]
  with $\gamma = 1/(180\zeta)$ and $\zeta=\mathcal{O}(1)$ encoding scheme-dependent normalization factors.
  In the natural spectral scheme one finds $\gamma=1/30$.
  This structure maps directly onto the parametrizations used in
  gravitational-wave catalogues and yields a bound on the pre-geometric
  scale $\ell_\chi$ from interferometric data.

  The Lorentzian completion thus embeds the spectral entropy dynamics
  in a fully causal framework and establishes gravitational waves as
  infrared probes of the underlying pre-geometric scale.
}
 
  \noindent\textbf{Keywords:} Spectral geometry, Emergent gravity, Lorentzian geometry, Gravitational waves,
  Modified dispersion relations, Effective field theory, Heat-kernel expansion, Seeley–DeWitt coefficients,
  Causal propagation, Quantum gravity phenomenology

  \clearpage
  \tableofcontents
  \newpage

  \section{Introduction}

  In previous work, we introduced a spectral entropy functional
  $S_\Pi[g] = \tfrac12 \log \det{}' A_g$,
  where $A_g$ is a Laplace-type elliptic operator associated with the
  projective substrate.
  In the infrared regime, the renormalized first variation of this functional
  was shown to generate an effective Einstein response, with higher-derivative,
  non-local, and anomaly contributions suppressed by powers of the
  pre-geometric scale $\ell_\chi$.
  This established the emergence of Einstein gravity as the dominant
  two-derivative sector of the spectral dynamics.

  However, the construction of Papers~1 and~2 was formulated in a
  Riemannian (elliptic) framework.
  While sufficient to derive the infrared field equations and their
  thermodynamic interpretation, this setting does not by itself address
  causal propagation.
  In particular, the elliptic operator $A_g$ encodes the spatial
  correlation structure of the substrate, but gravitational waves are
  dynamical, hyperbolic excitations propagating along null cones.

  The purpose of the present work is therefore to provide the
  Lorentzian completion of the spectral entropy dynamics.
  The Lorentzian character of this completion — specifically the
  signature $(-,+,+,+)$ and the four-dimensional structure of the
  effective spacetime — is not an additional assumption of the framework.
  Both are derived from the admissibility constraints of the Cosmochrony
  framework in the companion paper Q5b~\cite{Beau2026q5b}: the
  four-dimensionality follows from the homogeneous dimension
  $D_{\mathrm{hom}} = 4$ of $\mathrm{Heis}_3(\mathbb{R})$
  (Bass--Guivarc'h theorem), and the Lorentzian signature from the
  Born--Infeld bounded-flux constraint.
  The present paper takes these two structural results as input and
  constructs the causal dynamics on the resulting background.
  We define the Lorentzian action as
  \[
    S_\Pi^{(L)}[g]
    =
    \frac12 \Re \log \det{}' \mathcal{D}_g,
  \]
  equivalently via a Schwinger--Keldysh prescription,
  so that the quadratic kernel is real and retarded.
  The inverse operator entering the second variation
  $\delta^2 S_\Pi^{(L)}$ is thus the retarded Green operator $G_R$,
  ensuring causal propagation.

  We show that, in the regime
  $R\ell_\chi^2 \ll 1$,
  the quadratic operator derived from
  $\delta^2 S_\Pi^{(L)}$
  propagates exactly two transverse-traceless tensor modes of helicity
  $\pm 2$.
  Additional poles generated by higher-derivative corrections appear at
  $k^2 \sim \ell_\chi^{-2}$ and therefore lie outside the infrared
  window $\omega \ll \ell_\chi^{-1}$ relevant to present observations.
  In this domain, the theory reproduces the polarization content of
  general relativity.

  Beyond the leading Einstein sector, the spectral expansion generates a
  Weyl-squared contribution whose coefficient is fixed by the
  Seeley--DeWitt coefficient $a_4$.
  Linearization about Minkowski spacetime then yields a modified
  dispersion relation of the form
  \[
    \omega^2
    =
    c^2 k^2
    -
    \gamma\, \ell_\chi^2 k^4
    +
    \mathcal{O}(k^6),
  \]
  where $\gamma = 1/(180\zeta)$ and
  $\zeta = \mathcal{O}(1)$ encodes scheme-dependent normalization
  factors.
  In the natural spectral identification of the renormalization scale
  with the cutoff, $\mu=\ell_\chi^{-1}$, one has $\zeta=1/6$ and thus
  $\gamma=1/30$.

  This structure leads to a direct observational prediction.
  Gravitational-wave catalogues constrain modified dispersion relations
  of the form $\omega^2 = c^2 k^2 + A_4 k^4$.
  Comparison with the spectral result yields a bound
  \[
    \ell_\chi
    \lesssim
    \sqrt{180\,\zeta\,A_4^{\mathrm obs}},
  \]
  which becomes
  $\ell_\chi \lesssim \sqrt{30\,A_4^{\mathrm obs}}$
  in the natural scheme.
  The Lorentzian completion therefore provides the first direct
  gravitational-wave constraint on the pre-geometric scale introduced in
  the spectral framework.

  The structure of the paper is as follows.
  Section~\ref{sec:lorentzian_completion} defines the causal Lorentzian
  spectral action and its second variation.
  Section~\ref{sec:ir_spectrum} establishes the infrared propagating
  degrees of freedom and proves the transverse-traceless proposition.
  Section~\ref{sec:dispersion} derives the modified dispersion relation
  and its connection to the Seeley--DeWitt coefficients.
  Section~\ref{sec:gw_constraints} discusses the confrontation with
  gravitational-wave observations and the resulting bounds on
  $\ell_\chi$.
   \section{Lorentzian Completion of the Spectral Entropy Functional}
  \label{sec:lorentzian_completion}

  \subsection{From Elliptic to Hyperbolic Dynamics}
    \label{subsec:from-elliptic-to-hyperbolic-dynamics}

    In the Riemannian formulation developed previously, the spectral entropy
    functional is defined by
    \[
      S_\Pi[g] = \frac12 \log \det{}' A_g,
    \]
    where $A_g$ is a Laplace-type elliptic operator.
    Its renormalized first variation generates an infrared Einstein response,
    with higher-derivative corrections governed by the Seeley--DeWitt
    coefficients.

    While this formulation captures the static and quasi-static structure
    of the emergent geometry, it does not address causal propagation.
    In particular, $A_g$ is elliptic and therefore encodes spatial
    correlation structure rather than hyperbolic time evolution.
    To describe gravitational waves and dynamical perturbations,
    a Lorentzian completion is required.

    We therefore introduce a hyperbolic operator $\mathcal{D}_g$ such that its spatial restriction on constant-time
    slices reduces to $A_g$.
    In a $3+1$ decomposition with lapse $N$ and shift $N^i$, the d'Alembertian reads
    \[
      \Box_g = -\frac{1}{N^2} (\partial_t - \mathcal{L}_{\vec N})^2 + \Delta_\gamma + \text{(connection terms)},
    \]
    where $\Delta_\gamma$ is the Laplacian on spatial slices.
    In the quasi-static regime $N \simeq 1$, $N^i \simeq 0$, and for slowly varying extrinsic curvature, one recovers
    \[
      \Box_g \;\longrightarrow\; -\partial_t^2 + \Delta_\gamma.
    \]
    The elliptic operator $A_g$ of the Riemannian theory is thus identified
    with the spatial block of the Lorentzian operator.

  \subsection{Definition of the Lorentzian Spectral Action}
    \label{subsec:definition-of-the-lorentzian-spectral-action}

    The Lorentzian completion of the spectral entropy functional is defined as
    \[
      S_\Pi^{(L)}[g] = \frac12 \Re \log \det{}' \mathcal{D}_g.
    \]
    Equivalently, one may formulate the theory using a Schwinger--Keldysh contour prescription,
    so that the inverse operator entering functional variations is the retarded Green operator $G_R$.

    This choice ensures:

    \begin{itemize}
      \item Reality of the quadratic kernel.
      \item Causal propagation, i.e.\ $G_R(x,y)=0$ for $x \notin J^+(y)$.
      \item Consistency with classical gravitational-wave dynamics.
    \end{itemize}

    The operator $\mathcal{D}_g$ is assumed to be globally hyperbolic and
    normally hyperbolic in the sense of Lorentzian geometry, so that
    retarded and advanced Green operators exist and are unique.

  \subsection{Second Variation and Quadratic Operator}
    \label{subsec:second-variation-and-quadratic-operator}

    The second variation of the Lorentzian spectral action takes the operatorial form
    \[
      \delta^2 S_\Pi^{(L)} = \frac12 \mathrm{Tr} \left( G_R \,\delta^2 \mathcal{D}_g \right)
      -
      \frac12 \mathrm{Tr} \left(
        G_R \,\delta \mathcal{D}_g\,
        G_R \,\delta \mathcal{D}_g
      \right).
    \]

    This expression is the hyperbolic analogue of the elliptic formula derived previously.
    The replacement $A_g^{-1} \to G_R$ is the essential modification ensuring causal structure.

    The existence and uniqueness of retarded Green operators for normally
    hyperbolic operators on globally hyperbolic spacetimes are standard
    results of Lorentzian analysis~\cite{Baer2015}.
    The use of a Schwinger--Keldysh prescription ensures a causal and
    real effective action in curved spacetime~\cite{Jordan1986,Calzetta2008}.

    Expanding about Minkowski spacetime,
    \[
      g_{\mu\nu} = \eta_{\mu\nu} + h_{\mu\nu},
    \]
    one obtains a quadratic operator of the schematic form
    \[
      \mathcal{O} = c_{\mathrm{EH}} \Box + \beta \Box^2 + \text{(non-local terms)} + \mathcal{O}(R\ell_\chi^2).
    \]

    In the infrared regime $R\ell_\chi^2 \ll 1$, the dominant local contributions are the Einstein term and the
    Weyl-squared (Bach) correction.
    Non-local contributions are suppressed by additional powers of
    $\ell_\chi^2$ and are analytic in $\Box$ at low frequency.

  \subsection{Infrared Consistency}
    \label{subsec:infrared-consistency}

    The structure above guarantees that:

    \begin{enumerate}
      \item The leading kinetic operator is hyperbolic and luminal.
      \item Higher-derivative corrections introduce an additional pole at $k^2 \sim \ell_\chi^{-2}$.
      \item In the domain $\omega \ll \ell_\chi^{-1}$,
      the additional excitation decouples,
      leaving a standard massless graviton sector.
    \end{enumerate}

    The Lorentzian completion therefore preserves the infrared Einstein
    limit while embedding it in a fully causal framework.
    Subsequent sections analyze the propagating degrees of freedom and
    derive the resulting modified dispersion relation.
   \section{Infrared Spectrum and Propagating Degrees of Freedom}
  \label{sec:ir_spectrum}

  \subsection{Linearization Around Minkowski Spacetime}
    \label{subsec:linearization-around-minkowski-spacetime}

    We consider perturbations about flat spacetime,
    \[
      g_{\mu\nu}
      =
      \eta_{\mu\nu}
      +
      h_{\mu\nu},
    \]
    and work to quadratic order in $h_{\mu\nu}$.
    In the infrared regime
    $R\ell_\chi^2 \ll 1$,
    the quadratic operator derived from
    $\delta^2 S_\Pi^{(L)}$
    takes the local form
    \[
      \mathcal{O}_{\mu\nu}^{\ \ \rho\sigma}
      =
      c_{\mathrm{EH}}\Box
      +
      \beta \Box^2
      +
      \mathcal{O}(\ell_\chi^4 \Box^3),
    \]
    where $c_{\mathrm{EH}}>0$ and $\beta<0$ in the natural spectral scheme.

    We impose harmonic gauge,
    \[
      \partial^\mu \bar h_{\mu\nu}
      =
      0,
      \qquad
      \bar h_{\mu\nu}
      =
      h_{\mu\nu}
      -
      \frac12 \eta_{\mu\nu} h,
    \]
    which eliminates gauge redundancies at linear order.

  \subsection{Scalar--Vector--Tensor Decomposition}
    \label{subsec:scalar--vector--tensor-decomposition}

    To identify the propagating modes, we perform a
    scalar--vector--tensor (SVT) decomposition with respect to spatial
    rotations.
    Writing spatial indices as $i,j=1,2,3$, the perturbation splits as
    \[
      h_{00},
      \qquad
      h_{0i} = \partial_i B + B_i,
      \qquad
      h_{ij}
      =
      2\psi\,\delta_{ij}
      +
      2\partial_i\partial_j E
      +
      \partial_i E_j
      +
      \partial_j E_i
      +
      h^{\mathrm{TT}}_{ij},
    \]
    with
    \[
      \partial_i B_i = 0,
      \qquad
      \partial_i E_i = 0,
      \qquad
      \partial_i h^{\mathrm{TT}}_{ij} = 0,
      \qquad
      h^{\mathrm{TT}}_{ii} = 0.
    \]

    In harmonic gauge, the scalar and vector components enter the quadratic
    action without independent kinetic terms in the regime
    $R\ell_\chi^2 \ll 1$.
    They are therefore non-dynamical and can be eliminated algebraically.

    The only sector acquiring a propagating kinetic operator is the
    transverse-traceless tensor sector $h^{\mathrm{TT}}_{ij}$.

  \subsection{Transverse-Traceless Sector}
    \label{subsec:transverse-traceless-sector}

    Restricting to $h^{\mathrm{TT}}_{\mu\nu}$,
    the quadratic operator reduces to
    \[
      \left(
        c_{\mathrm{EH}}\Box
        +
        \beta \Box^2
      \right)
      h^{\mathrm{TT}}_{\mu\nu}
      =
      0.
    \]

    Factoring gives
    \[
      \Box
      \left(
        1
        +
        \frac{\beta}{c_{\mathrm{EH}}}
        \Box
      \right)
      h^{\mathrm{TT}}_{\mu\nu}
      =
      0.
    \]

    The first factor corresponds to the standard massless graviton pole
    $k^2=0$.
    The second factor introduces an additional pole at
    \[
      k^2
      =
      -\frac{c_{\mathrm{EH}}}{\beta}
      \sim
      \ell_\chi^{-2}.
    \]

    Since $\ell_\chi^{-1}$ sets the pre-geometric ultraviolet scale,
    this excitation lies outside the infrared window
    \[
      \omega \ll \ell_\chi^{-1},
    \]
    relevant to gravitational-wave observations.

  \subsection{Infrared Propagation Theorem}
    \label{subsec:infrared-propagation-theorem}

    \paragraph{Proposition.}
      In the regime $R\ell_\chi^2 \ll 1$ and $\omega \ll \ell_\chi^{-1}$,
      the Lorentzian spectral action propagates exactly two transverse-traceless tensor modes of helicity $\pm 2$.
      No scalar or vector propagating degree of freedom is dynamically accessible in this domain.

    \paragraph{Proof.}
      The SVT decomposition shows that scalar and vector components lack
      independent kinetic operators in harmonic gauge.
      The tensor sector obeys a fourth-order equation that factorizes into two second-order operators.
      The additional pole is located at $k^2 \sim \ell_\chi^{-2}$, hence outside the infrared regime.
      Only the massless pole contributes to low-frequency propagation.
      \hfill $\square$

\subsection{Polarization Content}
  \label{subsec:polarization-content}

  The surviving infrared excitations are precisely the two helicity $\pm 2$ tensor modes of general relativity.
  No scalar (breathing or longitudinal) or vector polarizations appear in the accessible frequency range.
  The Lorentzian completion of the spectral framework is therefore
  consistent with current gravitational-wave polarization constraints.

  \subsection{Sourced Vector Sector and the Gravitomagnetic Constraint}
    \label{subsec:sourced-vector-sector}

    The propagation theorem classifies the vector sector $h_{0i}$ as non-dynamical in the infrared.
    This is a statement about free propagation in vacuum, not a statement that the sector vanishes:
    an algebraically eliminated sector becomes a \emph{constraint} sector as soon as a source is present.
    We make this explicit by coupling the linearized operator to a projected stress tensor $T^{(\Pi)}_{\mu\nu}$,
    and we keep every conclusion below conditional on its momentum-density component $T^{(\Pi)}_{0i}$, whose
    microscopic construction from moving projected configurations belongs to the matter sector and is not
    attempted here.

    \paragraph{Sourced field equation.}
      Adding the minimal coupling $\half\!\int h_{\mu\nu}\,T^{(\Pi)\,\mu\nu}\,\dd^4 x$ to the quadratic action and
      extremizing in harmonic gauge gives, with the Einstein normalization
      $c_{\mathrm{EH}} = (16\pi G_{\mathrm{N}})^{-1}$ inherited from the infrared Einstein response,
      \[
        \Box\,\bar h_{\mu\nu} = -16\pi G_{\mathrm{N}}\, T^{(\Pi)}_{\mu\nu}.
      \]
      The $0i$ components isolate the vector sector,
      \[
        \boxed{\;\Box\,\bar h_{0i} = -16\pi G_{\mathrm{N}}\, T^{(\Pi)}_{0i}\;}
      \]
      which carries no second-order time derivative once the harmonic constraint is imposed.

    \paragraph{Stationary reduction.}
      Frame dragging is a stationary effect, sourced by a slowly rotating mass current, so $\partial_t \to 0$ and
      $\Box \to \nabla^2$.
      Writing the gravitomagnetic potential $A^{\mathrm{g}}_i \equiv -\tfrac14\,\bar h_{0i}$ and the mass current
      $j^{\mathrm{g}}_i \equiv T^{(\Pi)}_{0i}$, the momentum density $\rho\,v_i$ for dust, one obtains the vector
      Poisson equation
      \[
        \nabla^2 A^{\mathrm{g}}_i = 4\pi G_{\mathrm{N}}\, j^{\mathrm{g}}_i,
      \]
      structurally identical to magnetostatics $\nabla^2 \mathbf{A} = -\mu_0\,\mathbf{j}$.
      Its Green-function solution is
      \[
        \mathbf{A}^{\mathrm{g}}(\mathbf{x})
        = -G_{\mathrm{N}} \int \frac{\mathbf{j}^{\mathrm{g}}(\mathbf{x}')}{|\mathbf{x}-\mathbf{x}'|}\,\dd^3 x',
        \qquad
        \mathbf{B}^{\mathrm{g}} = \nabla \times \mathbf{A}^{\mathrm{g}},
      \]
      and for a source of angular momentum $\mathbf{J}$ the multipole expansion yields the dipolar
      Lense--Thirring potential and the associated precession.
      This last identification is a reproduction of linearized general relativity and carries no discriminating
      power: the framework reproduces general relativity in its domain of validity, so the measured frame dragging
      around the Earth is consistent with, but does not single out, the present construction.

    \paragraph{Consistency with the propagation theorem.}
      In the stationary regime the boxed equation is elliptic, a constraint, not a hyperbolic propagation law.
      This is precisely the content of the statement that the vector sector is non-dynamical in the infrared:
      the propagation theorem is used here, not bypassed.
      It guarantees the absence of a free vector mode, so that $\bar h_{0i}$ is fixed entirely by the source.
      The two helicity-$\pm2$ tensor modes propagate, while the vector sector is a sourced constraint, and the
      gravitomagnetic field is therefore determined without any additional degree of freedom.

    \paragraph{Inherited spectral correction.}
      Retaining the subleading operator $\beta\Box^2$ of the quadratic action, the stationary vector constraint
      becomes
      \[
        \Bigl(\nabla^2 + \tfrac{\beta}{c_{\mathrm{EH}}}\,\nabla^4\Bigr)\bar h_{0i}
        = -16\pi G_{\mathrm{N}}\, T^{(\Pi)}_{0i},
        \qquad
        \frac{\beta}{c_{\mathrm{EH}}} = -\gamma\,\ell_\chi^2,
        \quad \gamma = \frac{1}{180\,\zeta},
      \]
      where $\beta/c_{\mathrm{EH}}$ is the same ratio of Seeley--DeWitt coefficients that fixes the $k^4$
      dispersion correction of the transverse-traceless sector (Section~\ref{sec:dispersion}).
      The gravitomagnetic potential thus inherits a short-range correction of range $\sim\ell_\chi$, twin to the
      dispersion correction of the propagating modes.
      To our knowledge this is the only place where the gravitomagnetic sector would carry a framework-specific
      imprint rather than reproduce general relativity, and it is therefore the natural candidate for a
      discriminating statement, to be handled with the usual caution about magnitude.
      Its status is structural, conditional on extending the vacuum analysis of
      Section~\ref{subsec:infrared-propagation-theorem} to the sourced vector sector, and it is not yet phenomenologically
      evaluated.

    \paragraph{Status and localization of the gap.}
      Conditional on a supplied $T^{(\Pi)}_{0i}$, the reduction above fixes the gravitomagnetic component $g_{0i}$
      exactly as in linearized general relativity: the downstream step is structurally closed and introduces no new
      hypothesis.
      The open part is upstream, namely the construction of $T^{(\Pi)}_{0i}$ from moving projected matter
      configurations, which is the projected momentum density of the matter sector and remains an open deliverable
      of the gravitational programme.
      The derivation therefore does not close that gap; it localizes it entirely upstream, showing that the
      gravitomagnetic sector has no intrinsic difficulty of its own, only a dependence on the projected matter
      current.
   \section{Modified Dispersion from Spectral Coefficients}
  \label{sec:dispersion}

  \subsection{Quadratic Operator and Pole Structure}
    \label{subsec:quadratic-operator-and-pole-structure}

    In the transverse-traceless sector established in Section~\ref{sec:ir_spectrum},
    the quadratic operator takes the factorized form
    \[
      \mathcal{O}_{\mathrm{TT}} = c_{\mathrm{EH}}\Box + \beta \Box^2 = c_{\mathrm{EH}}
      \Box
      \left( 1 + \frac{\beta}{c_{\mathrm{EH}}}\Box \right).
    \]

    The coefficient $\beta$ originates from the Weyl-squared term in the renormalized spectral expansion.
    For a Laplace-type scalar minimal operator in four dimensions, the Seeley--DeWitt coefficient $a_4$ contains
    \[
      \alpha_C = \frac{1}{360(4\pi)^2},
    \]
    multiplying $C_{\mu\nu\rho\sigma}^2$.
    Using
    \[
      \delta \int \sqrt g\, C^2 = -2 \int \sqrt g\, B_{\mu\nu}\,\delta g^{\mu\nu},
    \]
    one obtains
    \[
      \beta = -2\alpha_C = -\frac{1}{180(4\pi)^2}.
    \]
    The heat-kernel expansion and the associated Seeley--DeWitt
    coefficients are reviewed in~\cite{Vassilevich2003,DeWitt1965}.
    The variation of the Weyl-squared term leading to the Bach tensor is
    standard in higher-derivative gravity~\cite{Bach1921,Stelle1977}.

    The Einstein prefactor is parameterized as
    \[
      c_{\mathrm{EH}} = \zeta\,(4\pi)^{-2}\,\ell_\chi^{-2},
    \]
    where $\zeta=\mathcal{O}(1)$ accounts for scheme-dependent
    normalizations and possible multiplicity factors.
    In the natural spectral identification of the renormalization scale
    with the cutoff, $\mu=\ell_\chi^{-1}$, and using $a_2=R/6$, one has $c_{\mathrm{EH}}=\mu^2/[6(4\pi)^2]$,
    corresponding to $\zeta=1/6$.

  \subsection{Infrared Dispersion Relation}
    \label{subsec:infrared-dispersion-relation}

    For plane-wave solutions $h^{\mathrm{TT}}_{\mu\nu} \sim e^{ik\cdot x}$, the operator yields the dispersion equation
    \[
      -c_{\mathrm{EH}} k^2 + \beta k^4 = 0.
    \]

    Solving for $\omega^2$ gives
    \[
      \omega^2 = c^2 k^2 - \frac{\beta}{c_{\mathrm{EH}}} k^4.
    \]

    Using the expressions above,
    \[
      \frac{\beta}{c_{\mathrm{EH}}} = -\frac{1}{180\,\zeta}\, \ell_\chi^2,
    \]
    so that the infrared dispersion relation becomes
    \[
      \omega^2 = c^2 k^2 - \gamma \ell_\chi^2 k^4 + \mathcal{O}(k^6),
    \]
    with
    \[
      \gamma = \frac{1}{180\,\zeta}.
    \]

    In the natural spectral scheme $\zeta=1/6$,
    \[
      \gamma = \frac{1}{30},
    \]
    and therefore
    \[
      \omega^2 = c^2 k^2 - \frac{1}{30}\, \ell_\chi^2 k^4 + \mathcal{O}(k^6).
    \]

  \subsection{Physical Interpretation}
    \label{subsec:physical-interpretation}

    The correction term is negative, implying a slight sub-luminal deviation of the group velocity at high frequency,
    \[
      v_g = \frac{\partial \omega}{\partial k} \simeq c \left( 1 - \frac{3}{2}\gamma \ell_\chi^2 k^2 \right).
    \]

    Since the modification is suppressed by $(k\ell_\chi)^2$, it is negligible in the regime $k\ell_\chi \ll 1$.
    The additional pole at $k^2 \sim \ell_\chi^{-2}$ corresponds to a heavy excitation beyond the
    infrared window probed by current interferometers.

    The modified dispersion relation is therefore fully consistent with the infrared Einstein limit,
    while providing a concrete $k^4$ correction that can be directly confronted with gravitational-wave observations.
   \section{Confrontation with Gravitational-Wave Observations}
  \label{sec:gw_constraints}

  \subsection{Mapping to the LVK Parametrization}
    \label{subsec:mapping-to-the-lvk-parametrization}

    Gravitational-wave catalogues constrain modified dispersion relations
    of the general form
    \[
      \omega^2 = c^2 k^2 + A_\alpha k^\alpha,
    \]
    with $\alpha=4$ corresponding to the leading higher-derivative
    correction compatible with locality and parity.

    The spectral Lorentzian completion derived above predicts
    \[
      \omega^2 = c^2 k^2 - \gamma \ell_\chi^2 k^4 + \mathcal{O}(k^6),
    \]
    with
    \[
      \gamma = \frac{1}{180\,\zeta}.
    \]

    Identifying coefficients gives
    \[
      A_4 = -\gamma \ell_\chi^2.
    \]

    Current analyses place bounds on $|A_4|$ from phase corrections
    accumulated during propagation over cosmological distances.
    Denoting the observational upper bound by $A_4^{\mathrm obs}$,
    the spectral prediction implies
    \[
      \ell_\chi \lesssim \sqrt{\frac{A_4^{\mathrm obs}}{\gamma}} = \sqrt{180\,\zeta\,A_4^{\mathrm obs}}.
    \]

    In the natural spectral scheme $\zeta=1/6$,
    this reduces to
    \[
      \ell_\chi \lesssim \sqrt{30\,A_4^{\mathrm obs}}.
    \]

  \subsection{Robustness and Scheme Dependence}
    \label{subsec:robustness-and-scheme-dependence}

    The coefficient $\zeta=\mathcal{O}(1)$ captures normalization
    conventions in the identification of the Einstein term and possible
    multiplicity factors associated with the underlying operator.
    While its precise numerical value is scheme-dependent,
    the following features are robust:

    \begin{itemize}
      \item The sign of the $k^4$ correction is negative.
      \item The structure of the dispersion is fixed by the ratio of Seeley--DeWitt coefficients.
      \item The constraint on $\ell_\chi$ scales as $\sqrt{\zeta}$ and is therefore stable up to factors of order unity.
    \end{itemize}

    Thus, gravitational-wave observations provide a direct infrared probe
    of the pre-geometric scale introduced in the spectral framework.

  \subsection{Luminal Front Propagation}
    \label{subsec:luminal-front-propagation}

    Multi-messenger observations constrain the difference between the
    speed of gravitational waves and the speed of light to extremely high precision.
    In the present framework, the characteristic operator remains
    $\Box$ at leading order, and the correction enters only at order $k^4$.
    The front velocity is therefore luminal, and deviations of the group
    velocity are suppressed by $(k\ell_\chi)^2$.

    Consequently, the Lorentzian spectral completion is consistent with
    current bounds on gravitational-wave speed while predicting a specific
    higher-order dispersive correction.

  \subsection{Summary of Observational Implications}
    \label{subsec:summary-of-observational-implications}

    Modified dispersion relations for gravitational waves are
    constrained by the LIGO--Virgo tests of general relativity on the
    GWTC-1 catalogue~\cite{LIGO2019Dispersion}.
    Multi-messenger observations further constrain deviations from
    luminal propagation~\cite{LIGO2017Speed}.

    The Lorentzian completion of the spectral entropy functional leads to:

    \begin{enumerate}
      \item Exactly two propagating tensor polarizations in the infrared.
      \item A modified dispersion relation with fixed structure $\omega^2 = c^2 k^2 - \gamma \ell_\chi^2 k^4$.
      \item A direct observational bound on the pre-geometric scale
      $\ell_\chi$ from gravitational-wave data.
    \end{enumerate}

    This establishes, for the first time within the spectral framework,
    a concrete and testable connection between the emergent geometric dynamics and interferometric observations.
   \section{Conclusion}
  \label{sec:conclusion}

  We have constructed the Lorentzian completion of the spectral entropy framework introduced in previous work.
  While the original formulation was elliptic and sufficient to derive
  the infrared Einstein response and its thermodynamic interpretation, it did not address causal propagation.
  The present paper establishes that the spectral dynamics admits a
  consistent hyperbolic completion with a well-defined retarded Green operator and a real quadratic kernel.

  Linearization about Minkowski spacetime shows that, in the regime
  $R\ell_\chi^2 \ll 1$ and $\omega \ll \ell_\chi^{-1}$,
  the theory propagates exactly two transverse-traceless tensor modes of helicity $\pm 2$.
  Scalar and vector sectors remain non-dynamical in the infrared,
  and the additional pole generated by higher-derivative corrections
  appears at $k^2 \sim \ell_\chi^{-2}$, beyond the frequency range of current gravitational-wave observations.
  The polarization content therefore coincides with that of general relativity in the accessible domain.

  Beyond the leading Einstein sector, the ratio of Seeley--DeWitt
  coefficients fixes a universal $k^4$ correction to the dispersion relation,
  \[
    \omega^2
    =
    c^2 k^2
    -
    \gamma \ell_\chi^2 k^4
    +
    \mathcal{O}(k^6),
  \]
  with $\gamma = 1/(180\zeta)$ and $\zeta=\mathcal{O}(1)$ encoding scheme-dependent normalization factors.
  In the natural spectral scheme, $\gamma=1/30$.
  This structure provides a direct mapping to the parametrizations used
  in gravitational-wave catalogues and yields a bound on the pre-geometric scale $\ell_\chi$ from observational data.

  The Lorentzian completion therefore closes the logical chain of the spectral program:
  the elliptic structure encodes spatial correlations,
  the renormalized first variation yields the infrared Einstein response,
  and the hyperbolic extension ensures causal propagation and observational testability.
  Gravitational waves thus provide a concrete infrared probe of the
  pre-geometric scale underlying the spectral construction.

  As a first step in this direction, we have shown that once a projected momentum density $T^{(\Pi)}_{0i}$ is
  supplied the vector sector obeys an elliptic gravitomagnetic constraint, fixing the frame-dragging potential
  $g_{0i}$ as in linearized general relativity, so that the remaining gap is localized entirely upstream in the
  construction of the projected matter current.
  Further developments include a fully covariant coupling to matter
  sectors and a systematic analysis of non-local corrections beyond the leading $k^4$ term.
  These directions may sharpen the connection between the spectral
  substrate and precision gravitational observations.
 
  \bibliographystyle{plain}
      \bibliography{cosmochrony-bibliography}

\end{document}
